\section{Aplicabilidade a edificações públicas universitárias}
\label{sec:aplicabilidade}

A caracterização documental identificou 93 registros sobre edificação genérica e 58 sobre portfólio predial. Entre as tipologias específicas, houve 17 registros sobre hospitais, 15 sobre edifícios comerciais, 12 residenciais, 11 universidades, dez \textit{campus}, cinco edifícios públicos, cinco escolas, quatro patrimônios históricos e um museu. As categorias são respostas múltiplas; edificação genérica é uma categoria-guarda-chuva, não uma tipologia homogênea.

\begin{table}[H]
\centering
\caption{Contextos de edificação identificados no núcleo final}
\label{tab:contexto}
\small
\begin{tabularx}{\textwidth}{Y >{\raggedleft\arraybackslash}p{2.1cm} >{\raggedleft\arraybackslash}p{2.1cm}}
\toprule
Contexto & Registros & \% dos 104 \\
\midrule
Edificação genérica & 93 & 89,4 \\
Portfólio predial & 58 & 55,8 \\
Hospital & 17 & 16,3 \\
Edifício comercial & 15 & 14,4 \\
Edifício residencial & 12 & 11,5 \\
Universidade & 11 & 10,6 \\
\textit{Campus} & 10 & 9,6 \\
Edifício público & 5 & 4,8 \\
Escola & 5 & 4,8 \\
Patrimônio histórico & 4 & 3,8 \\
Museu & 1 & 1,0 \\
\bottomrule
\end{tabularx}
\fonteautor
\end{table}

Em instituições de ensino, a gestão sustentável combina desempenho físico, usuários e capacidade administrativa. A literatura identifica fatores críticos em universidades \parencite{awuzie_fmuniversidades_2017}, critérios ambientais, econômicos e sociais em edifícios educacionais \parencite{alfalah_frameworkeducacional_2022}, indicadores operacionais de \textit{campus} \parencite{naji_campusdelphi_2023}, previsão orçamentária e uso de ordens de serviço \parencite{martins_custosmanutencaoses_2024,morais_eficienciamanutencao_2023}. O produto deve, portanto, registrar fonte e data dos dados, ausências, unidade comparada, vetos e histórico de pesos, usando entradas simples e exportáveis para prestação de contas.

A baixa frequência de universidade, \textit{campus} e edifício público confirma a lacuna de validação em portfólios heterogêneos, com restrições orçamentárias, continuidade acadêmica e múltiplos usuários. Em texto completo, a manutenção sustentável foi associada à resiliência institucional e aos ODS 4, 9 e 11 \parencite{ndukuba_resilienciainstitucional_2025}; um instrumento estrangeiro exigiu adaptação e concordância profissional em hospital público \parencite{talib_hospitalfm_2013}; e três projetos mostraram perda de informação na transição entre construção e operação, enfrentada pela integração de BIM, \textit{Soft Landings} e comunicação \parencite{tan_fluxoinformacao_2018}. Esses achados sustentam governança da informação e validação contextual, sem equivaler a validação do protocolo proposto.
